%% 阶段报告（IEEEtran 中文版）— 截至 2026-09-12 已完成部分
%% 编译：xelatex -interaction=nonstopmode stage_report_ieee_zh.tex （跑两遍）
%% 与 stage_report_ieee_en.tex 结构一致，两者独立可编译
\documentclass[conference]{IEEEtran}

\usepackage{amsmath}
\usepackage{amssymb}
\usepackage{booktabs}
\usepackage{graphicx}
\usepackage{float}
\usepackage{xeCJK}
\setCJKmainfont{SimSun}[BoldFont=SimHei,ItalicFont=KaiTi]
\setCJKsansfont{SimHei}
\setCJKmonofont{FangSong}
\usepackage{url}

\begin{document}

\title{DeReFusion 复现与金融时间序列算子专业化诊断\\ \large ——阶段性技术报告}

\author{\IEEEauthorblockN{Dongxu Jiang}
\IEEEauthorblockA{Email: jiangdongxu04@gmail.com}}

\maketitle

\begin{abstract}
本报告汇总截至 2026-09-12 已完成的工作：DeReFusion 的复现，以及围绕同一个问题的一串受控诊断——
输入的结构状态是否决定哪种算子更有效。先复现论文核心对比（DLinear 线性基座、LSTM--Transformer
残差分支、无参数加法融合，统一置于 RevIN），随后按因果相对波动率对预测误差分层。非线性分支的边际
价值是\textbf{资产依赖}而非普适：标普 500 上显著为正且跨种子稳健，ETHUSD 方向一致但统计临界，
BTCUSD 显著\textbf{反号}。容量敏感性实验（隐层宽度 $23\!\to\!128$，其余冻结）表明早先"线性全面
占优"的读数\textbf{部分来自非线性算子的容量瓶颈}：容量充分时 ETHUSD 在 $36/36$ 个观测上偏好非线性，
而 BTCUSD 从不偏好。容量对齐的代理路由实验（12 个预先固定的结构状态、线性 vs.\ MLP、协议一致）
未发现稳健的状态级算子专业化。最后，十个资产的探索性扫描显示交互效应\textbf{跨零}（六负四正，
五个显著），并识别出两个 LOO 稳定的关联（已实现波动率与 $|\mathrm{ACF}_1|$）——仅作为
\textbf{候选}资产级结构规律报告。报告明确区分：证据排除了什么（样本级路由）、支持了什么
（资产级算子异质性）、以及仍未解决的是什么。
\end{abstract}

\begin{IEEEkeywords}
复现；时间序列预测；算子专业化；结构状态；自适应融合；受控对比。
\end{IEEEkeywords}

\section{研究范围与问题}
本阶段遵循"证据优先"流程：先复现参照结果，再一次只问一个问题，且所有分组规则、阈值与算子定义都在
看到结果\textbf{之前}固定。三个问题：\textbf{Q-A} 已发表的融合策略排序能否在独立数据上复现；
\textbf{Q-B} 非线性分支的边际贡献是否取决于输入窗口的波动状态；\textbf{Q-C} 输入的\textbf{结构状态}
是否决定哪种算子更有效（路由是否有实验依据）。全部实验在 CPU 上运行，采用冻结的 70/10/20 时序切分，
报告逐样本配对 bootstrap 置信区间（4000 次）与配对胜率；观测到测试结果后未调任何超参数。

\section{参照对比的复现}
上游数据集不可再分发，故重新抓取十个资产的日线 OHLC（2016--2025）；行数与论文一致（股指与个股 2513，
外汇 2602，BTCUSD 3653，ETHUSD 2975）。表~\ref{tab:repro} 给出标普 500 在 $T\in\{1,24\}$ 的对比
（seed 2021）。

\begin{table}[!t]
\caption{核心对比的复现（GSPC，seed 2021，归一化尺度）}
\label{tab:repro}
\centering
\small
\begin{tabular}{@{}lcccc@{}}
\toprule
模型 & $T{=}1$ & $T{=}24$ & $R^2$ & 参数量 \\
\midrule
DeReFusion（加法）     & 0.01647 & \textbf{0.06230} & \textbf{0.8691} & 28{,}436 \\
revin-DLinear（线性）  & \textbf{0.01577} & 0.07007 & 0.8528 & 4{,}664 \\
gatev1（波动率门控）   & ---     & 0.06767 & 0.8578 & 28{,}580 \\
gatev2（可学习门控）   & 0.02526 & 0.07191 & 0.8489 & $\sim$28k \\
\bottomrule
\end{tabular}
\end{table}

$T=24$ 上论文排序复现：加法融合的 MSE 比线性基座好 $11.1\%$，两个门控变体都\textbf{劣于}无参数加法。
$T=1$ 上线性基座略优，与长线性预测器的"先分解后简单模型"行为一致。扩展到另外两个资产（$T=24$）：
BTCUSD $0.21797$ 对 $0.22322$（$+2.4\%$）、ETHUSD $0.14602$ 对 $0.14876$（$+1.8\%$），
加法优势排序为 GSPC $\gg$ BTCUSD $\approx$ ETHUSD。

\section{波动率工况分层}
每个测试窗口的工况由\textbf{仅用输入窗口}的因果波动率度量确定：对数收益标准差及去趋势版本
$RV^{rel}_t = RV_t / \mathrm{median}(RV_{t-120:t-1})$。表~\ref{tab:regime} 报告非线性分支相对线性
基座的分层改善（主口径为相对波动率）与交互效应 $\Delta_{high50}-\Delta_{low50}$。

\begin{table}[!t]
\caption{分层表现（相对波动率）。$\Delta=$ MSE$_{非线性}-$MSE$_{线性}$，负值偏向非线性。
$^\dagger$ CI 不含 0。}
\label{tab:regime}
\centering
\small
\begin{tabular}{@{}lcccc@{}}
\toprule
资产 & 最平静 20\% & 高波动半区 & 最剧烈 20\% & 交互效应 \\
\midrule
GSPC (2021) & $+6.0\%$ & $+14.2\%^\dagger$ & $+21.6\%^\dagger$ & $-0.01205^\dagger$ \\
GSPC (2022) & $-2.7\%$ & $+23.2\%^\dagger$ & $+29.6\%^\dagger$ & $-0.02248^\dagger$ \\
ETHUSD      & $+16.7\%^\dagger$ & $+6.3\%^\dagger$ & $+14.2\%^\dagger$ & $-0.00943$ \\
BTCUSD      & $+3.5\%$ & $-1.0\%$ & $-3.5\%$ & $+0.01502^\dagger$ \\
\bottomrule
\end{tabular}
\end{table}

两点结论。其一，标普 500 的效应\textbf{跨种子稳健}：两个种子都给出同量级且显著为负的交互效应，
平静工况下两分支不可分。其二，BTCUSD 显著反号，因此"非线性优势随波动率上升"**不能**作为普适判据。
方法学检查解释了此前部分困惑：原始波动率与时间强混淆（$RV$ 与样本序号相关：GSPC $+0.44$、BTCUSD
$-0.67$）；去趋势后的 $RV^{rel}$ 消除该趋势，其下 GSPC 中间五分位的"线性反超"消失。
报告全部以 $RV^{rel}$ 为主口径。

\section{自适应融合没有收益}
对工况结构的自然利用是门控 $\alpha(X)$ 混合两分支，框架本身恰好提供该变体。在结构最强、跨种子最稳健
的资产上运行（GSPC，$T=24$）：波动率感知门控 MSE $0.06767$，比无参数加法（$0.06230$）\textbf{差
$8.6\%$}，仅比纯线性基座（$0.07007$）好 $3.4\%$。由于非线性分支在平静工况从不劣于线性分支，
最优门控权重实际上恒为 1，门控\textbf{没有可切换的空间}。结合此前可学习门控的负结果，证据支持
\textbf{停止自适应融合这条线}，而非在其上继续搭路由器。

\section{容量敏感性}
路由代理实验用容量对齐的 MLP 对比线性算子（$9{,}240$ 对 $9{,}431$ 参数，$+2.1\%$），
该对齐把 384 维输入压进 23 维隐层——严重瓶颈。因此把隐层宽度作为\textbf{唯一变量}
（$23\to64\to128$；数据、特征、切分、优化器、学习率、批大小、轮数、早停、种子与评估协议全部不变）
重跑结构状态评估。表~\ref{tab:cap} 给出 12 状态 $\times$ 3 种子的平均 $\Delta$MSE。

\begin{table}[!t]
\caption{容量敏感性。平均 $\Delta$MSE（负 $=$ 非线性更优）；括号内为 36 个（状态,种子）观测中偏好
非线性的个数。}
\label{tab:cap}
\centering
\small
\begin{tabular}{@{}lcccc@{}}
\toprule
隐层 & 参数量 & GSPC & BTCUSD & ETHUSD \\
\midrule
23  & 9{,}431  & $+0.0598$ (0/36) & $+0.1783$ (0/36) & $+0.0145$ (12/36) \\
64  & 26{,}200 & $+0.0250$ (7/36) & $+0.1888$ (0/36) & $\mathbf{-0.0281}$ (36/36) \\
128 & 52{,}376 & $+0.0061$ (8/36) & $+0.2998$ (0/36) & $\mathbf{-0.0321}$ (36/36) \\
\bottomrule
\end{tabular}
\end{table}

三项预先指定的检查一致。（i）\textbf{符号反转}：宽度间共 25 例状态级反转，23 例属 ETHUSD（三种子皆有），
8 例属 GSPC（**全部且仅 seed 2023**，即种子不稳定）；BTCUSD 零例。（ii）\textbf{状态间跨度}：
GSPC（$0.106\to0.060$）与 ETHUSD（$0.089\to0.040$）的 $\Delta$MSE 状态间极差随容量\textbf{缩小}；
BTCUSD 的扩大来自单一高误差状态的量级放大。（iii）\textbf{逐种子方向一致性}由 $31/36$ 降至 $29/36$、
$28/36$。结论：容量充分后 ETHUSD 对非线性算子的偏好是\textbf{资产级全体}（十二个状态全数反转，而非
子集）——这是\textbf{资产级异质性}而非\textbf{状态级专业化}。据此，早先"线性全面占优"记为
\textbf{部分由容量造成}，且不把容量扫描当作路由证据。

\section{结构感知路由：最小受控代理实验}
\label{sec:routing}
\subsection{设计}
十二个结构状态在观察任何结果\textbf{之前}定义，使用因果窗口特征：滞后 1/5/10 自相关、趋势 $t$ 统计量、
符号持续性、阈值固定为 $c=3\sigma$ 的跳变比例、去趋势相对波动率、偏度与峰度。二值状态以训练集中位数
为阈值；另报告复合的\textbf{结构强度}状态与 $K{=}3$ 的 $k$-means 交叉检验。两个算子为线性映射
（$9{,}240$ 参数）与单隐层 23 单元的 MLP（$9{,}431$ 参数），协议完全一致（同种子、优化器、步数、
批大小、学习率、早停与预处理）。

\subsection{结果}
容量对齐下，GSPC 与 BTCUSD 在\textbf{全部}十二个状态上线性算子更优（GSPC $\Delta$MSE 介于
$+0.047$ 与 $+0.081$，三种子全显著；BTCUSD 介于 $+0.129$ 与 $+0.448$）。只有 ETHUSD 出现两种符号的
状态间差异（结构弱状态偏向非线性，最多 $-0.013$；结构强状态偏向线性，$+0.033$）。没有状态的偏好能
通过容量扫描检验。\textbf{主要发现：容量对齐尺度上未建立稳健的算子专业化}；唯一跨资产不变式是
平静工况下两算子不可分。

\section{十个资产的资产级异质性}
用同一冻结协议（$T=24$，seed 2021）在另外七个资产上运行，以刻画交互效应的\textbf{符号}。
GSPC 的两个种子合并为一个资产，故 $N=10$。表~\ref{tab:assets} 报告交互效应，
表~\ref{tab:assoc} 报告资产结构特征与该效应的探索性关联（含 leave-one-asset-out，LOO 检查）。

\begin{table}[!t]
\caption{资产级交互效应（相对波动率）。$^\dagger$ CI 不含 0。}
\label{tab:assets}
\centering
\small
\begin{tabular}{@{}lccc@{}}
\toprule
资产 & $\Delta_{interaction}$ & 95\% CI & 显著 \\
\midrule
EURUSD & $-0.1371$ & $[-0.164,-0.111]$ & 是$^\dagger$ \\
GSPC   & $-0.0173$ & $[-0.027,-0.009]$ & 是$^\dagger$ \\
USDJPY & $-0.0122$ & $[-0.026,+0.002]$ & 否 \\
ETHUSD & $-0.0094$ & $[-0.019,+0.000]$ & 否 \\
DJI    & $-0.0084$ & $[-0.019,+0.003]$ & 否 \\
TM     & $-0.0081$ & $[-0.028,+0.011]$ & 否 \\
SOX    & $+0.0044$ & $[-0.008,+0.017]$ & 否 \\
BTCUSD & $+0.0150$ & $[+0.006,+0.024]$ & 是$^\dagger$ \\
BABA   & $+0.0405$ & $[+0.022,+0.060]$ & 是$^\dagger$ \\
NVO    & $+0.0976$ & $[+0.046,+0.149]$ & 是$^\dagger$ \\
\bottomrule
\end{tabular}
\end{table}

\begin{table}[!t]
\caption{资产结构特征与 $\Delta_{interaction}$ 的探索性 Spearman 关联（$N=10$）与 LOO 稳健性。}
\label{tab:assoc}
\centering
\small
\begin{tabular}{@{}lcccc@{}}
\toprule
特征 & $\rho$ & $p$ & LOO 范围 & 稳健 \\
\midrule
$|\mathrm{ACF}_1|$  & $+0.733$ & $0.016$ & $[+0.633,+0.817]$ & 是 \\
已实现波动率        & $+0.661$ & $0.038$ & $[+0.533,+0.900]$ & 是 \\
峰度                & $+0.345$ & $0.328$ & $[+0.100,+0.533]$ & 是 \\
趋势持续性          & $+0.042$ & $0.907$ & $[-0.317,+0.283]$ & 否（变号） \\
跳变比例            & ---      & ---     & --- & 恒定 \\
\bottomrule
\end{tabular}
\end{table}

效应符号\textbf{跨零}：六负四正，五个统计显著。两个特征 LOO 稳定：已实现波动率更高、一阶自相关更强
的资产，在湍流工况下的非线性优势\textbf{更弱}（甚至反转）。在 $N=10$ 下两个关联均达 $p<0.05$
（$|\mathrm{ACF}_1|$：$\rho=+0.733$，$p=0.016$；已实现波动率：$\rho=+0.661$，$p=0.038$），
SOX 的加入（交互 $+0.0044$，统计不定）不改变判定。
\textbf{这些关联是探索性的，仅作为候选规律报告}：十资产下单个 $p$ 值不足以定论；跳变比例特征在所有资产的中位数上恒定（每 95 个收益中
1 次超阈），在预注册阈值下不带区分信息；趋势持续性在 LOO 下变号，标记为单资产敏感。另有两点限定：
四个 CI 含零的资产（USDJPY、ETHUSD、DJI、TM）方向全部为负、与显著负例一致，负号一侧内部一致；
两个显著正例（NVO、BABA）恰是一阶自相关最强的资产，与候选规律的预测一致——关联并非由显著性
格局偶然驱动，而是由效应量在特征取值范围内的排序驱动。

\section{诊断与决策}
\begin{table}[!t]
\caption{五个诊断问题的汇总回答}
\label{tab:q}
\centering
\small
\begin{tabular}{@{}p{0.46\columnwidth}p{0.48\columnwidth}@{}}
\toprule
问题 & 回答 \\
\midrule
Q1 是否存在稳健的样本级路由证据？ &
否。代理实验与容量扫描都未给出稳定的"状态 $\rightarrow$ 算子"映射。 \\
Q2 是否存在真实的资产级算子异质性？ &
是。容量充分时 ETHUSD 在 36/36 个观测上偏好非线性；BTCUSD 在任何容量下偏好线性；GSPC 接近打平。 \\
Q3 能否用现有结构特征解释？ &
仅能作为候选：已实现波动率与 $|\mathrm{ACF}_1|$ 与交互效应符号相关且通过 LOO，
但符号在十资产上跨零。 \\
Q4 当前主要不确定性？ &
资产依赖与算子形态——不是结构表示，也不再是路由（路由已基本排除）。 \\
Q5 是否可重新考虑更强的非线性算子？ &
仅可作为候选之一，须与容量对齐的 MLP 对照，且不得由"ETHUSD 偏好非线性"外推到任一算子族。 \\
\bottomrule
\end{tabular}
\end{table}

本阶段记录的决策为：\textbf{停止自适应融合}——波动率门控实测劣于无参数加法，且它想利用的结构本身
没有可切换的空间；\textbf{不构建神经路由器}——状态到算子的映射未被建立；把算子选择当作
\textbf{资产级配置问题}（这是证据支持的做法）；任何结构化非线性算子都须以"可重复的环境"与
"容量对齐的对比"为前提。

\section{局限}
\textbf{算子保真度}：代理算子（展平窗口上的线性映射与 MLP）不同于框架内分支——后者把 DLinear 分解
与 LSTM--Transformer 残差在 RevIN 之下组合。代理结果\textbf{不能}读作"非线性无用"；框架内残差分支
在标普 500 的 $T=24$ 上仍跨种子优于线性基座，两者并不矛盾。\textbf{样本量}：资产级关联建立在十个
资产、每资产一个种子（GSPC 两个）之上，关联系数的置信区间很宽，分析为探索性；四个 CI 含零的资产
（USDJPY、ETHUSD、DJI、TM）方向一致但只能描述为\textbf{方向性为负、统计上不定}。
\textbf{工况定义}：工况为因果波动率的五分位，平静工况在算子间统计不可分，因此"平静"是接受性结论而非独立效应。
\textbf{环境}：全部为 CPU 运行，墙钟时间与论文 GPU 数值不可比，仅解释相对比较。

\section{结论}
复现在 $T=24$ 上重现了已发表排序，且跨种子稳健。其后的诊断细化但不推翻该结果：非线性分支的边际价值
是\textbf{资产依赖}而非工况普适；波动率门控没有收益；容量对齐的对比表明，早先"线性一律占优"的偏好
\textbf{部分是容量瓶颈造成的假象}。我们\textbf{未}发现支持样本级路由的稳健实验依据，而确实发现了
真实的资产级算子异质性。SOX 的加入未改变 PATH 判定；在 $N=10$ 下两个 LOO 稳定特征的 $p$ 值均低于
0.05（$|\mathrm{ACF}_1|$：$p=0.016$；已实现波动率：$p=0.038$），但仍只作为\textbf{候选}规律报告。
因此证据支持的下一步问题不是"路由器应该如何选择"，而是
\textbf{"在哪些可重复的环境中，更强的非线性算子值得占用它的容量"}——且必须以容量对齐的基线为参照来回答，
并以逐资产的配置决策为落地形式，而不是按样本的动态路由。

\appendix[可复现性说明]
全部复现资产位于项目仓库的 \texttt{reproduction/} 下：分析脚本在 \texttt{reproduction/analysis/}，
批次运行器在 \texttt{reproduction/batches/}，各设置结果在 \texttt{reproduction/results/}。
四条环境事实曾耗费实际时间：数据管线无条件导入 \texttt{patool}、\texttt{huggingface\_hub}、
\texttt{sktime} 与 \texttt{datasets}；\texttt{joblib} 须 pin 到 \texttt{1.5.3}；CPU 运行必须加
\texttt{--no\_use\_gpu}；\textbf{杀掉 DataLoader 工作进程（\texttt{spawn\_main}）会让父训练进程死锁}，
"卡死"必须靠\textbf{进程 CPU 时间}而非日志输出来判断。

\begin{thebibliography}{1}
\bibitem{hsieh2026} Hsieh, C.-Y., Chen, Y.-C. (2026). DeReFusion: a decomposition-reconstruction
fusion framework for financial time series forecasting. \emph{Applied Soft Computing}, 203, 116252.
\bibitem{zeng2023} Zeng, A., Chen, M., Zhang, L., Xu, Q. (2023). Are transformers effective for
time series forecasting? \emph{Proc. AAAI}, 37(9), 11121--11128.
\end{thebibliography}

\end{document}